\section{Pennies with an external sink}\label{sec:attachments}
Over a field of characteristic two, consider vertices $x,y$, a loop $r$ at $x$, and arrows $p:x\to y$, $q:y\to x$, with
\begin{equation}\label{eq:two-penny}
 pq=r^2+\lambda r^3,\qquad qp=0,\qquad r^4=0,
 \qquad \lambda\in\{0,1\}.
\end{equation}
The three attachments below have the additional vertices, arrows and relations displayed in the table. Their outside vertices are sinks. The defining ideal is generated by the displayed relations and \eqref{eq:two-penny}.
\[
\begin{array}{c|c|l}
 \text{type}&\text{additional arrows}&\text{additional relations}\\\hline
 a&\beta:x\to b&r\beta=0\\
 b&\gamma:y\to c&rp\gamma=0\\
 c&\beta:x\to b,\ \gamma:y\to c&r\beta=q\beta=p\gamma=0
\end{array}
\]
We also use type $a$ with $q\beta=0$ and type $b$ with $p\gamma=0$. Denote each resulting algebra by $A_\lambda$, retaining its type. The standard choice is $\lambda=0$. In every type, $r^2p=qr^2=0$ follows from \eqref{eq:two-penny}: multiply $pq$ by $p$, or by $q$ on the left, and use that $1+\lambda r$ is invertible.

\subsection{The two exceptional modules}
For types $a,c$ and the smaller type $a$, take $J=kr^3$, $P=e_xA_\lambda$, and $I=P/S_b$. For type $b$ and its smaller quotient, put $j=qrp$, $J=kj$, $P=e_yA_\lambda$, and $I=P/S_c$. Each $J$ is a two-sided ideal. In the first case the spaces of $P$ have bases
\[
 P_x=(1,r,r^2,r^3),\quad P_y=(p,rp),\quad P_b=(\beta);
\]
its $c$-space in type $c$ is zero. In the second case they have bases
\[
 P_y=(e_y,j),\quad P_x=(q,qr),\quad P_c=(\gamma).
\]
In both cases $I$ is obtained by deleting the indicated sink vector.

\begin{lemma}\label{lem:exception-exhaustion}
For these algebras, $I$ is injective. The indecomposable modules not annihilated by $J$ are exactly $P$ and $I$.
\end{lemma}
\begin{proof}
First identify $I$ with the dual of the left projective at the same vertex as $P$. In types $a,c$, pair $I$ with $A_\lambda e_x$ by taking the coefficient of $r^3$ in the product. The root powers pair nondegenerately by complementary degrees summing to three. Between the right basis $p,rp$ and the left basis $q,qr$, the matrix is
\[
 \begin{pmatrix}\lambda&1\\1&0\end{pmatrix}.
\]
In type $b$, use the coefficient of $j$. The pairs $(e_y,j)$ with $(e_y,j)$ and $(q,qr)$ with $(p,rp)$ both have matrix $\left(\begin{smallmatrix}0&1\\1&0\end{smallmatrix}\right)$. Products at different vertices do not mix these blocks. Each pairing is nondegenerate and satisfies $\langle va,u\rangle=\langle v,au\rangle$. Hence it gives a right-module isomorphism $I\simeq D_k(A_\lambda e_x)$ or $D_k(A_\lambda e_y)$, respectively. Thus $I$ is injective with essential simple socle at its generating vertex. The extra relations of the two smaller types leave these bases and pairings unchanged.

Write $s$ for that generating vertex and $t$ for the deleted sink. There is an exact sequence
\begin{equation}\label{eq:exception-sequence}
 0\longrightarrow S_t\longrightarrow P\xrightarrow{\pi}I\longrightarrow0.
\end{equation}
The copy $PJ\simeq S_s$ in $P$ maps isomorphically to $\soc I$. If an indecomposable $M$ satisfies $MJ\ne0$, a vector on which the generating path of $J$ acts nontrivially defines a map $f:P\to M$ nonzero on $PJ$. Let $K=\ker f$. If $\pi(K)\ne0$, essentiality implies that it contains $\soc I$. Lift a nonzero socle vector to $K$; it differs from the corresponding vector in $PJ$ only by a vector at $t$. Multiplication by $e_s$, with $s\ne t$, removes the latter. This contradicts injectivity of $f$ on $PJ$. Therefore $K\subseteq S_t$.

The map $\operatorname{im}f\to I$ induced by $\pi$ extends, by injectivity of $I$, to a surjection $M\to I$. Let its kernel be $N$. Applying $\Hom(-,N)$ to \eqref{eq:exception-sequence}, using projectivity of $P$, gives
\begin{equation}\label{eq:exception-ext}
 \Ext^1(I,N)=N_t/\im(N_s\to N_t)= (\topm N)_t.
\end{equation}
The last equality uses the unique arrow entering the sink. Since $S_t$ is projective, write $N=N_0\oplus S_t^d$, with $(\topm N_0)_t=0$. Then $\Ext^1(I,N_0)=0$ and $\Ext^1(I,S_t^d)=k^d$. The zero column gives middle module $I\oplus N$, and any nonzero column can be taken by $\operatorname{GL}_d(k)$ to the first coordinate, giving $P\oplus N_0\oplus S_t^{d-1}$. Indecomposability forces $M\simeq I$ or $P$.
\end{proof}

\subsection{Comparison for common-quotient endpoints}
For types $a,c$, the common quotient is $C=A_\lambda/(r^3)$ in the displayed arrow coordinates. Put $R_I=\{v\in I:vr^3=0\}$ and $R_P=\{v\in P:vr^3=0\}$. Then $R_P=R_I\oplus S_b$. Identify their two versions by \eqref{eq:radical-identification} on the root and the identity elsewhere. The modules $T_P=P/PJ=e_xC$ and $T_I=I/IJ=T_P/S_b$ are common. Their natural boundaries differ by the root shear $1+r$. For $E=P^u\oplus I^v$, with $u,v\ge0$,
\[
 \rad(E/EJ)=H^{u+v}\oplus S_b^u,
\]
where $H$ is the two-vertex $H$ of Section~\ref{sec:boundary}. All its outside arrows act by zero.

\begin{lemma}\label{lem:attachment-comparison}
For each of the five displayed models, identify the common-quotient modules and replace $P,I$ by their standard counterparts. If $X,Y,N$ are common-quotient modules, then
\[
 F^{P^u\oplus I^v\oplus N}_{X,Y}(A_1)
 =F^{(P^0)^u\oplus(I^0)^v\oplus N}_{X,Y}(A_0)
 \qquad(u,v\ge0).
\]
This equality holds after every scalar extension.
\end{lemma}
\begin{proof}
In types $a,c$, repeat Lemma~\ref{lem:isolated-compensation} for $\Omega\subseteq H^{u+v}\oplus S_b^u$. The additional sink space is sent to zero in each correction. The root radical of the target is $N'_xr+N'_yq$. In the projective factorization, $r^2\beta=qr\beta=0$; in type $c$ the arm value $qrp$ is killed by $\gamma$, since $p\gamma=0$. The top-to-socle correction is a module map because its root image is killed by $\beta$ as well as $r,p$. For the smaller type $a$, $q\beta=0$ adds no further condition to these maps. Thus \eqref{eq:stable-compensation} holds in the actual common quotient, and Lemma~\ref{lem:graph-comparison} applies.

For type $b$, take instead $C=A_\lambda/(j)$. In $A_1/(j)$ the arrow $p'=(1+r)p$ satisfies
\[
 p'q=r^2,\qquad qp'=0,\qquad rp'\gamma=0.
\]
It identifies this quotient with $A_0/(j)$, where the standard $p$ corresponds to $p'$. In these common coordinates the annihilators of $j$ in $P,I$ are identified by
\[
 q\longmapsto q+qr,\qquad qr\longmapsto qr,\qquad j\longmapsto j,
\]
and by the identity at the sink. The common $T_P=P/PJ=e_yC$ and $T_I=I/IJ=T_P/S_c$ have root chain $q\mapsto qr\mapsto0$. Their $y$-space is generated by $e_y$, which $q$ sends to $q$; the common $p$ acts by zero. Only $T_P$ has the sink vector. The two boundaries differ on the root by $q\mapsto qr$ and agree elsewhere. Thus
\[
 (\rad T_E)_x\simeq (k[r]/(r^2))^{u+v},\qquad
 (\rad T_E)_y=0,\qquad (\rad T_E)_c=k^u,
\]
and all arrows between distinct vertices act by zero there.

For $\Omega\subseteq\rad T_E$, use the same top-to-socle construction as in Lemma~\ref{lem:isolated-compensation}. It reduces the remaining difference to a root map sending a chosen length-two-chain generator to $ur$, killing the bottom and all length-one chains, where
\[
 u=wr+zq\in(\rad N')_x,\qquad ur^2=up=0.
\]
Here $qr^2=0$, so $wr^3=0$. In $e_xC$, $r^3$ is nonzero. The first projective factor must therefore send the source length-two chain to $r^2,r^3$, and its length-one summands to zero. This commutes with $r$ and $p$ by $r^4=r^2p=r^3p=0$. Sending the projective generator to $w$ gives the required $wr^2$ on the upper layer and zero on the lower layer. The second factor sends the upper layer to $qr\in e_yC$ and the other components to zero; it is a module map because $qr^2=qrp=0$ in $C$. Sending $e_y$ to $z$ gives $zqr$. These constructions are linear in all generator coefficients and apply to arbitrary multiplicities. There is no condition from $\gamma$, since the source $y$-space is zero. They also work in the smaller quotient $p\gamma=0$.

The root socle correction is killed by both $r$ and $p$, so it is again an endomorphism of the entire module. This proves \eqref{eq:stable-compensation} for type $b$. The shear preserves every source submodule, and Lemma~\ref{lem:graph-comparison} gives the equality. All constructions use linear algebra over the base field and repeat after scalar extension.
\end{proof}

\subsection{A recursion for arbitrary endpoints}
We give the endpoint argument separately. Suppose $A$ is representation-finite, $C$ is a quotient, and the only indecomposables outside $\operatorname{mod}C$ are $P,I$. Assume that $P$ is projective, $I$ is injective, $\rad P\in\operatorname{mod}C$, and \eqref{eq:exception-sequence} holds with $\topm P=S_s$, where $t\ne s$ is a sink with a unique incoming arrow from $s$. Work first with a split model. Suppose Hall polynomials in $\ZZ[T]$ are known whenever both endpoints are $C$-modules, allowing an arbitrary middle module.

For any modules $U,V$, use the automorphism polynomials \eqref{eq:aut} and put
\begin{equation}\label{eq:split-coefficient}
 c(U,V)=\frac{a_{U\oplus V}(T)}{a_U(T)a_V(T)T^{h(U,V)}}\in\R^\times.
\end{equation}
It is the Hall coefficient of the split middle module. To justify that description, a nonsplit sequence $0\to V\to Z\to U\to0$ satisfies
\begin{equation}\label{eq:end-strict}
 \dim\End Z<\dim\End(U\oplus V).
\end{equation}
Indeed, left exactness gives
\[
 h(Z,Z)\le h(U,Z)+h(V,Z),\quad
 h(U,Z)\le h(U,V)+h(U,U)-1,\quad
 h(V,Z)\le h(V,V)+h(V,U).
\]
The strict term occurs because $\Hom(U,Z)\to\End U$ cannot contain the identity without splitting the sequence. Thus the only extension class with middle module $U\oplus V$ is zero, and \eqref{eq:extension-hall} gives \eqref{eq:split-coefficient}.

\begin{proposition}\label{prop:exception-recursion}
Under these hypotheses, $A$ has Hall polynomials.
\end{proposition}
\begin{proof}
First compute the entire row with quotient $I^m$. Write $N=N_0\oplus S_t^d$, where $N_0$ has no $S_t$ summand. Formula~\eqref{eq:exception-ext} identifies $\Ext^1(I^m,N)$ with $d\times m$ matrices. Rank $r$ has middle module
\[
 L_r=N_0\oplus P^r\oplus I^{m-r}\oplus S_t^{d-r},
 \qquad 0\le r\le\min(m,d).
\]
Row and column operations give this decomposition from \eqref{eq:exception-sequence}, and the $P$ multiplicity distinguishes the different ranks. Factoring a rank-$r$ matrix as an injection times a surjection and dividing by $\operatorname{GL}_r$ gives its count
\[
 \prod_{j=0}^{r-1}\frac{(Q^m-Q^j)(Q^d-Q^j)}{Q^r-Q^j}.
\]
For $r=0$ the product is one. Therefore
\begin{equation}\label{eq:injective-row}
 F^{L_r}_{I^m,N}=
 \frac{a_{L_r}(T)}{a_{I^m}(T)a_N(T)T^{h(I^m,N)}}
 \prod_{j=0}^{r-1}\frac{(T^m-T^j)(T^d-T^j)}{T^r-T^j}
 \in\R,
\end{equation}
and all other middle coefficients in this row are zero.

Induct on the dimension $\dim_kM$ of the middle module. Zero endpoints give the usual Kronecker coefficients. At a conservative extension the indecomposable labels form the same finite, nonrepeating list, as proved in Section~\ref{sec:defects}. Thus every bounded-dimension sum below has a fixed finite index set, and Hom and End dimensions remain unchanged.

If the quotient contains $P$, projectivity implies that the middle module also contains $P$. For $M=P\oplus M'$ the exact split-factor identity is
\begin{equation}\label{eq:projective-strip}
 F^{P\oplus M'}_{P\oplus X',Y}
 =\frac{c(P,M')}{c(P,X')}F^{M'}_{X',Y}.
\end{equation}
It follows from \eqref{eq:extension-hall} and
$h(P,M')=h(P,X')+h(P,Y)$. If the required summand is absent, the coefficient is zero. Similarly injectivity gives
\begin{equation}\label{eq:injective-strip}
 F^{M'\oplus I}_{X,Y'\oplus I}
 =\frac{c(M',I)}{c(Y',I)}F^{M'}_{X,Y'}.
\end{equation}
Both right sides have smaller middle dimension.

For three nonzero modules $U,V,W$ of total dimension $\dim_kM$, associativity gives
\begin{equation}\label{eq:hall-assoc}
 \sum_Z F^Z_{U,V}F^M_{Z,W}
 =\sum_N F^N_{V,W}F^M_{U,N}.
\end{equation}
It follows by counting two-step submodule flags. Every inner coefficient has middle dimension less than $\dim_kM$, so it is already in $\R$.

First let $Y\in\operatorname{mod}C$. The cases $X\in\operatorname{mod}C$, $X$ containing $P$, and $X=I$ are known. Otherwise $X=I\oplus X'$ with $X'\ne0$. For fixed $\dim X$, proceed in increasing order of $\dim\End X$. Formula~\eqref{eq:hall-assoc} gives
\begin{equation}\label{eq:recursion-one}
 c(I,X')F^M_{X,Y}
 =\sum_N F^N_{X',Y}F^M_{I,N}
   -\sum_{Z\not\simeq X}F^Z_{I,X'}F^M_{Z,Y}.
\end{equation}
The first sum uses the known row \eqref{eq:injective-row}; every nonzero term in the second has $\dim\End Z<\dim\End X$ by \eqref{eq:end-strict}. This constructs all columns with $Y$ a $C$-module. Formula~\eqref{eq:injective-strip} then constructs those with $Y=I^a\oplus Y_0$, $Y_0\in\operatorname{mod}C$.

Next consider $Y=P$. Its unique maximal submodule is $\rad P$, so $F^P_{S_s,\rad P}=1$. For $X\ne0$, associativity gives
\begin{equation}\label{eq:recursion-two}
 F^M_{X,P}
 =\sum_N F^N_{X,S_s}F^M_{N,\rad P}
  -\sum_{Z\not\simeq P}F^Z_{S_s,\rad P}F^M_{X,Z}.
\end{equation}
The first sum uses $C$-submodule columns. In the second, $\dim Z=\dim P$, so $Z\not\simeq P$ cannot contain $P$ and is a sum of $I$ and $C$-modules. These columns are already known. The inner dimensions are strictly smaller than $\dim_kM$, since $X\ne0$ and $\rad P\ne0$.

Finally, submodule endpoints without $P$ or with an $I$ summand have been treated. For the remaining $Y=Y'\oplus P$, $Y'\ne0$, use increasing End dimension at fixed $\dim Y$ and
\begin{equation}\label{eq:recursion-three}
 c(Y',P)F^M_{X,Y}
 =\sum_N F^N_{X,Y'}F^M_{N,P}
  -\sum_{Z\not\simeq Y}F^Z_{Y',P}F^M_{X,Z}.
\end{equation}
The first sum uses \eqref{eq:recursion-two}, and the second has strictly smaller End dimension by \eqref{eq:end-strict}. Fixed dimension has finitely many labels, so both inner orderings terminate. All divisions are by units \eqref{eq:split-coefficient}. These formulas construct a single element of $\R$ for every coefficient, equal to the actual Hall number on every conservative extension. Such field sizes are unbounded, for example with extension degrees congruent to one modulo $D_A$. Lemma~\ref{lem:integer} makes each function an element of $\ZZ[T]$.
\end{proof}

\begin{corollary}\label{cor:attachments}
Every representation-finite algebra among the five split models $A_1$ above has Hall polynomials.
\end{corollary}
\begin{proof}
For each model, Lemma~\ref{lem:exception-exhaustion} identifies its indecomposables with the common-quotient labels together with $P,I$. The same argument applies to $A_0$, so its finite representation type follows from that of $C$. The standard choice is ray-standard and hence geometrically standard by Lemma~\ref{lem:ray-mesh}. Lemma~\ref{lem:attachment-comparison} and Theorem~\ref{thm:standard} supply the $C$-endpoint input for Proposition~\ref{prop:exception-recursion}. The residues of $P,I$ are the base field, already represented by the common simple at their generating vertex. Consequently $D_{A_1}=D_C=D_{A_0}$, and the input holds on every required conservative extension. The radical of $P$ is annihilated by $J$, and \eqref{eq:exception-sequence} has its asserted unique sink arrow. Thus all the recursion hypotheses hold for each model itself.
\end{proof}
